\subsection{Algebras over Maximal Ordered Fields}

Let R be a maximal ordered field (VI, §2, No.~5, p.~25). Let C be the quadratic R-algebra of type \((-1,0)\) ; if \((1,i)\) is its canonical basis, then we have \(i^{2}=-1\) . Moreover, C is an algebraic closure of R (VI, §2, No.~6, p.~26, Theorem 3). Let H be the quaternion R-algebra of type \((-1,0,-1)\) . The multiplication table of H in its canonical basis \((1,i,j,k)\) is given by

\[
i ^ {2} = j ^ {2} = k ^ {2} = - 1, \quad i j = - j i = k, \quad - i k = k i = j, \quad j k = - k j = i.
\]

We identify C with the subalgebra \(R + Ri\) of H. The conjugate of an element \(q = x + yi + zj + tk\) of H is \(\overline{q} = x - yi - zj - tk\) . The Cayley trace and norm of q are given by

\[
\mathrm{T} (q) = q + \overline {{q}} = 2 x, \quad \mathrm{N} (q) = q \overline {{q}} = x ^ {2} + y ^ {2} + z ^ {2} + t ^ {2}.
\]

Since R is an ordered field, we have \(N(q) > 0\) if \(q \neq 0\) , so H is a field, with center R (VIII, p.~363, Proposition 2). The reduced trace and norm of an element q of H are \(T(q)\) and \(N(q)\) , respectively.

THEOREM 1. --- Let D be an R-algebra of finite degree that is a field. Then D is isomorphic to R, C, or H.

Denote the center of D by Z, and let L be a maximal commutative subfield of D. We have \([D:Z]=[L:Z]^{2}\) by VIII, p.~265, Corollary 2; we also have \([L:R]\leqslant2\) because C is an algebraic closure of R. There are consequently three possible cases:

\begin{enumerate}
\def\labelenumi{\alph{enumi})}
\item
  We have \(\mathrm{R} = \mathrm{Z} = \mathrm{L}\) , so \([\mathrm{D}:\mathrm{Z}] = 1\) and \(\mathrm{D} = \mathrm{R}\) .
\item
  We have \(\mathrm{R} \neq \mathrm{Z}\) and \(\mathrm{Z} = \mathrm{L}\) , so \([\mathrm{D} : \mathrm{Z}] = 1\) and \(\mathrm{D} = \mathrm{L}\) . In this case, \(\mathrm{D}\) is isomorphic to \(\mathrm{C}\) .
\item
  We have R = Z and \([L : R] = 2\) , so \([D : R] = 4\) . By Proposition 4 of VIII, p.~364, the R-algebra D is isomorphic to a quaternion algebra of type \((\alpha, 0, \gamma)\) , where \(\alpha\) and \(\gamma\) are nonzero elements of R. Let \(i \in D - Z\) be such that \(i^{2} = \alpha\) . We have \(\alpha \neq 0\) . If \(\alpha > 0\) , then there exists an \(a \in R\) such that \(a^{2} = \alpha\) (VI, §2, No.~6, p.~26, Theorem 3); we then have \((a - i)(a + i) = 0\) , which is absurd because D is a field. So we have \(\alpha < 0\) . The inequality \(\gamma < 0\) is shown analogously. There then exist elements a and c of \(R^{*}\) such that \(\alpha = -a^{2}\) and \(\gamma = -c^{2}\) (loc. cit.). The algebra D is therefore isomorphic to the quaternion algebra of type \((-1,0,-1)\) (VIII, p.~362, Remark 2), that is, to H.
\end{enumerate}

Remarks. --- 1) Let O be the octonion algebra of type \((-1,0,-1,-1)\) over R (III, Appendix, No.~3, p.~615). Let D be an alternative Cayley algebra over R such that every nonzero element of D has an inverse. We can prove (VIII, p.~370, Exercise 5) that D is isomorphic to R, C, H, or O.

*2) The above applies to the field \(\mathbf{R}\) of real numbers. Every \(\mathbf{R}\) -algebra of finite degree that is a field is isomorphic to \(\mathbf{R}\) , \(\mathbf{C}\) , or \(\mathbf{H}\) .

\begin{enumerate}
\def\labelenumi{\arabic{enumi})}
\setcounter{enumi}{2}
\tightlist
\item
  Let A be a normed algebra over the field R. Suppose that A is a field. Then A is isomorphic to R, C, or H (``Gelfand--Mazur theorem'') (cf.~Comm. Alg., VI, §6, No.~4, p.~407, Theorem 1 and TS, I, §2, n° 5, p.~26, corollaire 2).*
\end{enumerate}

